\documentclass[11pt]{article}
\usepackage{amsmath,amssymb,amsthm}
\usepackage[margin=1in]{geometry}
\newtheorem{proposition}{Proposition}
\title{Variance-derivative cancellation in the discounted moment closure}
\author{Research note accompanying the numerical comparison}
\date{14 November 2026}
\begin{document}
\maketitle
\noindent\textbf{Status.} This note proves an algebraic identity for smooth
fields. It does not prove local existence of the nonlinear calibration system.
The numerical comparison is documented in \texttt{benchmarks/RESULTS.md}.

\section{Definitions}
Let $z=(x,u,v)$, where $x=\log(S/S_0)$, and write
\[
 s(t,x)=\sigma(t,S_0e^x),\quad a=s^2>0,\quad
 g=\partial_x s=\frac{a_x}{2\sqrt a},\quad
 \beta=\rho_{Sd}\nu_d,\quad\gamma=\rho_{Sf}\nu_f,
 \quad c=\rho_{df}\nu_d\nu_f.
\]
The physical rates are $r_d=f_d^0+\eta_d+u$ and
$r_f=f_f^0+\eta_f+v$, with
$\eta_i=\nu_i^2 B_{\kappa_i}(t)^2/2$ and
$B_k(t)=(1-e^{-kt})/k$ (continuously extended at $k=0$).
In particular $r_d,r_f$ have no explicit $x$ dependence.
Set
\[
 b=(r_d-r_f-a/2,-\kappa_d u,-\kappa_f v-\gamma s),\qquad
 Q=\begin{pmatrix}a&\beta s&\gamma s\\
 \beta s&\nu_d^2&c\\\gamma s&c&\nu_f^2\end{pmatrix},
 \qquad h=(s,\beta,\gamma)^\top.
\]
The killed forward operator is
\[
 \mathcal K_s^*m=-\nabla\cdot(bm)
       +\frac12\sum_{i,j}\partial_{ij}(Q_{ij}m)-r_dm.
\]
The five-field system evolves the discounted physical density $p$ and signed
moments $(m_1,m_2,m_3,m_4)$ of
$(\ell_d,\zeta_f,\ell_f,\zeta_d)$. Its two $\ell$ equations are
\begin{align}
 \partial_t m_1&=\mathcal K_s^*m_1-(\kappa_d+sg)m_1-m_2
 +(V_d+\beta s)p-\nabla\cdot(ghm_1),\\
 \partial_t m_3&=\mathcal K_s^*m_3-\kappa_fm_3+m_4
 +(-V_f+\gamma s)p-\nabla\cdot(ghm_3),
\end{align}
where $V_i(t)=\nu_i^2 B_{2\kappa_i}(t)$.

\section{Exact cancellation}
Define the non-divergence operators
\begin{align*}
 \mathcal D_a={}&\frac a2\partial_{xx}
 +\frac{\nu_d^2}2\partial_{uu}+\frac{\nu_f^2}2\partial_{vv}
 +\beta s\partial_{xu}+\gamma s\partial_{xv}+c\partial_{uv},\\
 \mathcal A_a={}&\mathcal D_a+
 \left(-r_d+r_f+\frac a2+\frac{a_x}2\right)\partial_x
 +\kappa_d u\partial_u+(\kappa_f v+\gamma s)\partial_v.
\end{align*}
\begin{proposition}
For positive $s$ twice continuously differentiable in $x$ and sufficiently
smooth moment fields, the two equations above are exactly equivalent to
\begin{align}
 \partial_t m_1&=\mathcal A_a m_1+(\kappa_f-r_d)m_1-m_2
                         +(V_d+\beta s)p,\label{eq:active1}\\
 \partial_t m_3&=\mathcal A_a m_3+
       \left(\kappa_d-r_d+\frac{a_x}2\right)m_3+m_4
                         +(-V_f+\gamma s)p.\label{eq:active3}
\end{align}
Neither equation contains $a_{xx}$.
\end{proposition}
\begin{proof}
Because $sg=a_x/2$, the pure log-spot terms satisfy
\[
 \tfrac12\partial_{xx}(am)-\partial_x(sgm)
       =\tfrac12a m_{xx}+\tfrac12a_xm_x.
\]
The mixed terms satisfy
\[
 \beta\partial_{xu}(sm)-\beta\partial_u(gm)=\beta s m_{xu},
 \qquad
 \gamma\partial_{xv}(sm)-\gamma\partial_v(gm)=\gamma s m_{xv}.
\]
Expanding the drift divergences therefore gives
\[
 \mathcal K_s^*m-\nabla\cdot(ghm)
 =\mathcal A_am+
 (\kappa_d+\kappa_f-r_d+a_x/2)m.
\]
Subtracting $(\kappa_d+sg)m_1$ in the domestic equation or $\kappa_fm_3$
in the foreign equation proves the claim, including the asymmetric remaining
$a_x/2$ term in \eqref{eq:active3}.
\end{proof}

\section{What it says about nonlinear calibration}
For $\bar m_i=\int_{\mathbb R^2}m_i\,du\,dv$ and
$n=\bar m_1-\bar m_3$, consider the known-market-denominator feedback
\[
 a=q-\frac{2n}{p_*},\qquad
 p_*(t,x)=S_0e^x C_{KK}(t,S_0e^x),
\]
where $q$ is the target Dupire variance. Equivalence with the model-denominator
feedback still requires the regularity and identification argument in the
separate discounted-closure note; it is not assumed for arbitrary fields.

For any of the passive fields $R=p,m_2,m_4$, expansion of
$\mathcal K_s^*R$ contains the term $R a_{xx}/2$. The other reaction terms
in those equations involve at most $a_x$. Consequently their highest feedback
derivative is
\[
 -\frac{R}{p_*}\partial_{xx}\int_{\mathbb R^2}(m_1-m_3)\,du\,dv.
\]
Derivatives of $1/p_*$ multiply only $n$ or $n_x$. By contrast,
\eqref{eq:active1}--\eqref{eq:active3} contain at most the first derivative of
this marginal feedback. Ordering the active fields $(m_1,m_3)$ first, and
$(p,m_2,m_4)$ second, thus gives a triangular structure to the highest
\emph{feedback} derivatives. Every row still contains the ordinary physical
second-order diffusion $\mathcal D_a$.

This identity removes a particular derivative coupling. It does not by itself
establish sectoriality, coercivity or an invertible linearised operator in a
function space adequate for the whole-space calibration problem.

\section{Remaining steps for an existence theorem}
A meaningful perturbative result would fix an admissible Gaussian reference
with $a_0\geq\underline a>0$ and prove that a nonzero neighbourhood of market
smile perturbations admits calibrated solutions, with stability in a stated
norm. The following issues remain unresolved in the present work:
\begin{enumerate}
\item A Gaussian marginal $p_*$ tends to zero in the wings. Positivity at every
finite strike does not bound $1/p_*$ on the whole line. A weighted space must
control the marginal projection, division, products and derivatives together.
\item The passive equations retain second derivatives of active marginals.
A full linearised estimate, including unbounded Gaussian-rate drifts, killing
and nonlocal couplings, is needed. Formal triangularity is insufficient.
\item The original initial law is a point mass. Starting a perturbation at
positive time after a Gaussian evolution avoids this initial singularity,
but would be a narrower theorem and must be labelled accordingly.
\item The nonlinear estimate must preserve the variance margin and identify
the resulting fields with a positive discounted density and physical moments.
\end{enumerate}
No claim of nonflat local existence, stability, or removal of regularisation
is made until these estimates are supplied. The flat Gaussian benchmark,
the exact linear closure, the cancellation, and the finite numerical
calibrations are distinct results.
\end{document}
